\chapter{DX-016 --- Error Surface for Crash Post-Mortems}\label{record-321}

\noindent{\footnotesize\color{BookMuted}\path{docs/method/graveyard/DX-016-error-surface-post-mortem.md}\par}\smallskip

\begin{RecordMeta}
title: "DX-016 — Error Surface for Crash Post-Mortems"
legend: DX
lane: graveyard

\end{RecordMeta}

\section{Disposition}\label{record-321-disposition}

Superseded by \href{https://github.com/flyingrobots/bijou/blob/cbd70612140e7a450d764b3c7be1c6af878b99bc/docs/method/retro/DX-021-error-surface-post-mortem.md}{DX-021 --- Error Surface for Crash Post-Mortems}. The crash-surface work shipped, but this older backlog note used the \texttt{D\allowbreak{}X\allowbreak{}-\allowbreak{}0\allowbreak{}1\allowbreak{}6} identifier before that id was later reused for scenario-provenance bench work. Keep this file only as historical lineage for the legacy backlog tree.

\section{Original Proposal}\label{record-321-original-proposal}

Legend: \href{https://github.com/flyingrobots/bijou/blob/1853ad902f6bf219326409e7f743a5a1d6517dbc/docs/method/legends/DX-developer-experience.md}{DX --- Developer Experience}

\section{Idea}\label{record-321-idea}

When a TUI app crashes during \texttt{u\allowbreak{}p\allowbreak{}d\allowbreak{}a\allowbreak{}t\allowbreak{}e\allowbreak{}(\allowbreak{})} or \texttt{v\allowbreak{}i\allowbreak{}e\allowbreak{}w\allowbreak{}(\allowbreak{})}, the \texttt{r\allowbreak{}u\allowbreak{}n\allowbreak{}t\allowbreak{}i\allowbreak{}m\allowbreak{}e\allowbreak{}.\allowbreak{}t\allowbreak{}s} currently catches the error, writes the stack trace to stderr, and exits. If the app is in alt-screen mode, this stack trace is often lost or obscured by the terminal's screen restoration.

Implement a dedicated ``Error Surface'' that is rendered using Bijou components \emph{before} the alt-screen is exited. This surface should display: - The error message and stack trace. - The last stable model snapshot (if serializable). - A prompt to ``Press Enter to exit.''

\section{Why}\label{record-321-why}

\begin{enumerate}
\def\labelenumi{\arabic{enumi}.}
\tightlist
\item
  \textbf{Visibility}: Keeps critical crash data on-screen for the developer to inspect before the process dies.
\item
  \textbf{Professionalism}: Moves the app from a ``noisy CLI crash'' to a ``structured post-mortem.''
\item
  \textbf{Diagnostic Power}: Capturing the model state at the point of failure is an industrial-strength debugging aid.
\end{enumerate}

\section{Effort}\label{record-321-effort}

Medium --- requires a special ``Crash Mode'' in the run loop that bypasses the normal TEA update and renders a static diagnostic surface.

